\section{Positive auxiliary caps and a finite native extension interface}
\label{sec:r11-native-capped-presentation}

\paragraph{Status, authorship, and scope.}
This R11 module is constructed jointly by the R11 root and cube
agents and is submitted for exact-source independent review.
The first two lemmas
are actual finite-presentation statements in every rank.  The final
interface is \textbf{CONDITIONAL}: it does not provide the required
universal quadruple, assert a contraction on arbitrary actual forms,
or solve the raw cube.  Its purpose is to make precise how a supplied
finite homogeneous formula could retain the original field, all
proper-entry zeros, all cap walls, and the original metric.

\subsection{A positive unit-log in place of a constant cap offset}

Use the exact CS notation for a rank-$r$ diagonal square, $r\ge3$,
on $W=N\setminus|f=0|$, and let $a,b$ be its unit valuations.
On the original-$K$ product
\[
 \widehat W=W\times_K\{T\in\mathbf G_m^{\rm an}:v(T)>0\},
 \qquad \tau=v(T),
 \tag{NCP1}
\]
pull back the supplied norm and metric data by their finite native
presentations.  Define
\[
 \Psi_\tau=\min\{\psi,\Phi_0+|a|+|b|+\tau\},\qquad
 R_\tau=R-D\Psi_\tau,
 \tag{NCP2}
\]
where $\psi$ is the passive norm potential, with value $+\infty$
at its zero set.  The formula is one specified finite piecewise
affine valuation presentation after the common metric shift.  It
does not invoke closure under minima of arbitrary actual scalars.

\begin{lemma}[The whole CS formula with a positive unit-log cap]
The full expression CS36, with $\Psi_\tau$ in place of its original
constant-offset cap, is an actual triple on all of $\widehat W$ and
satisfies
\[
 P_s=\partial_2q_s[\Psi_\tau]+e_{a,b}R^{s-1}.
 \tag{NCP3}
\]
This is an identity with all parameter, cap-wall, metric, and mixed
coefficients included.
\end{lemma}

\begin{proof}
At a passive-zero point the ratio of passive to full norms is zero,
whereas $\exp(-|a|-|b|-\tau)$ is positive and continuous.  Thus the
finite branch in (NCP2) holds on an ordinary neighborhood, so the
cap is actual there.  On the passive-nonzero locus both entries of
the minimum have the indicated finite native presentation.  These
formulas agree on their ordinary overlaps.

The two ordinary opens
\[
 \{\psi<\Phi_0+|a|+|b|+\tau\},\qquad
 \{\psi>\Phi_0+|a|+|b|\}
 \tag{NCP4}
\]
cover $\widehat W$ because $\tau>0$.  On the first the normalizer
is exactly $\psi$, so the frozen CS21 passive calculation applies,
with the same genuine projective map and all base directions.
On the second every passive entry is strictly inactive in all four
norms by CS28.  The exact CS17--CS20 two-entry calculation allows
an arbitrary common reference potential.  It therefore applies to
the full $\Psi_\tau$ and $R_\tau$ there.  In particular no
$d'\tau$, $d''\tau$, or mixed cap-curvature coefficient is discarded.
The full residual identity CS29 follows by actual locality.  CS30--CS36
then prove (NCP3), using the complete MT formula and only the
independently closed pair $e_{a,b}$ for CPT.
\end{proof}

For any $\varpi\in K$ with $0<|\varpi|<1$, the original-$K$
constant section $T=\varpi$ specializes these finite formulas to
the legal constant cap $\gamma=v(\varpi)>0$.  This is substitution
in the stated native scalar/differential presentation, with
$d'\tau=d''\tau=0$ afterwards; it is not a higher-form descent
from a field extension.  If the six original cube faces have
different positive constants $\varepsilon_\nu$, the already proved
CT theorem compares each face, with its same passive block, from
$\gamma$ to $\varepsilon_\nu$.  Explicitly,
\[
 \mathcal Q_{\boldsymbol\varepsilon}
       -\mathcal Q_\gamma
   =\partial_3\sum_{\nu=1}^6\sigma_\nu
             \mathcal T^{\rm CT}_{\nu;\gamma,\varepsilon_\nu}.
 \tag{NCP5}
\]
No comparison of different passive blocks is asserted by CT.

\subsection{Simultaneous norm-preserving replacement near zero entries}

The following local lemma is deliberately stated for a finite
collection of supplied norms and supplied capped norms, not for
all analytic functions or all actual forms.

Let $f=(f_1,\ldots,f_r)$ be analytic functions in a local frame
near $x\in W$.  Supply finitely many full norms
\[
 M_\nu=\max_j|u_{\nu j}f_j|,
 \tag{NCP6}
\]
where every $u_{\nu j}$ is an analytic unit.  Also supply finitely
many nonempty passive subsets $I_\mu$, unit coefficients
$w_{\mu j}$, and passive norms
\[
 P_\mu=\max_{j\in I_\mu}|w_{\mu j}f_j|.
 \tag{NCP7}
\]
The supplied capped logarithms have the form
\[
 \Xi_\mu=\min\{-\log P_\mu+\varrho,
                  -\log M_{\sigma(\mu)}+\varrho+h_\mu\},
 \tag{NCP8}
\]
where $\varrho$ is the common local metric term and each $h_\mu$
is a finite continuous supplied actual scalar, unchanged in the
replacement.  The minima themselves are assumed supplied by the
finite native presentation, as in (NCP2).

\begin{lemma}[One ordinary neighborhood preserving the entire finite data]
There exist constants $\epsilon_j\in K$ and an ordinary neighborhood
$O$ of $x$ such that
\[
 g_j=f_j+\epsilon_j\quad\hbox{are analytic units on }O,
 \tag{NCP9}
\]
and on the \emph{entire} $O$ every full norm (NCP6), every capped
logarithm (NCP8), and the common metric term are unchanged by
$f\mapsto g$.  One may take $\epsilon_j=0$ whenever $f_j(x)\ne0$.
Thus all their actual derivatives, products, and mixed jets agree.
\end{lemma}

\begin{proof}
At $x$ some $f_j$ is nonzero, and hence every full norm $M_\nu(x)$
is positive.  Let $J$ be the finite set of indices with $f_j(x)=0$.
Choose nonzero $\epsilon_j\in K$ for $j\in J$, simultaneously so
small that
\[
 |u_{\nu j}(x)\epsilon_j|<M_\nu(x)
       \quad\hbox{for all }\nu,j\in J.
 \tag{NCP10}
\]
For every passive norm with $P_\mu(x)>0$, also require
\[
 |w_{\mu j}(x)\epsilon_j|<P_\mu(x)
                   \quad(j\in I_\mu\cap J).
 \tag{NCP11}
\]
If $P_\mu(x)=0$, all indices in $I_\mu$ belong to $J$; require
\[
 |w_{\mu j}(x)\epsilon_j|
   <M_{\sigma(\mu)}(x)e^{-h_\mu(x)}
                   \quad(j\in I_\mu).
 \tag{NCP12}
\]
There are finitely many strictly positive upper bounds.  Arbitrarily
small nonzero powers of the original-$K$ element $\varpi$ satisfy
them; no residue-field rationality of $x$ is needed.

All unchanged nonzero entries remain units on a neighborhood.
Each changed entry has value $\epsilon_j\ne0$ in $\mathscr H(x)$,
so it too is a unit after shrinking.  Shrink once more so that
(NCP10)--(NCP12), with their corresponding continuous norm bounds,
hold throughout one common ordinary neighborhood.  In (NCP10), an
unchanged entry attains the original full maximum near $x$ and the
modified zero entries remain strictly below that maximum.  Thus
every full norm is literally unchanged there.  The same argument
applies to each initially nonzero passive norm using (NCP11).

For an initially zero passive norm, shrink also so that the original
passive entries satisfy the bound on the right of (NCP12).  Both the
original and modified passive maxima are then strictly below the
same cap floor.  Their two capped logarithms therefore equal the
same second branch of (NCP8).  Hence every displayed equality is
an equality of functions on one ordinary open, not only at $x$.
Actual locality and the actual differential operations give all
the stated mixed-jet equalities.
\end{proof}

The lemma applies at once to the eight full diagonal cube norms
and its six finite passive caps, including caps with the extra
positive unit-log coordinate.  It does \emph{not} preserve an
uncapped passive logarithm that was infinite at $x$.  It does not
replace the original tuple or original cycle in a Green problem:
the temporary $g$ is used solely to check an identity between the
already specified original finite scalar formulas on a neighborhood.

\subsection{The precise conditional interface for a homogeneous formula}

Let $\mathscr P=\mathbf P^{r-1}_K$ with its standard metric on
$\mathcal O(1)$ and curvature $R_{\rm st}$.  On
\[
 \mathscr P^{\rm an}\times(\mathbf G_m^3)^{\rm an}
             \times\{v(T)>0\},
 \tag{NCP13}
\]
use the homogeneous coordinate tuple, the three distinct diagonal
unit changes, and all six caps of (NCP2) with the same $\tau$.
Write $\theta_v$ for its eight norm potentials; $\theta_0=0$
for the standard metric.  The six capped potentials are finite
global actual scalars.  On each projective ratio torus all these
functions are homogeneous piecewise linear in the ratio valuations,
the three unit valuations, and $\tau$.  Thus all their walls are
central in that enlarged coordinate space.  This last assertion
would be false if one kept a nonzero numerical cap offset while
omitting the extra coordinate.

Suppose a concrete finite quadruple $\mathscr V$ is supplied,
written entirely using real linear combinations, products and
$d',d''$ of the following global data:
\[
 R_{\rm st},\quad \theta_v,\quad
 \Psi_\nu,\quad a,b,c,\tau.
 \tag{NCP14}
\]
There must be no individual uncapped coordinate logarithm, division
by a possibly zero coordinate, unspecified primitive, or contraction
operation left in the final formula.  Suppose its identity is proved
in the complete native presentation calculus, not only on chamber
interiors or against a restricted test family:
\[
 \mathcal Q_\tau
       =\partial_3\mathscr V+R_{\rm st}^{r-3}E_{a,b,c},
 \quad
 E_{a,b,c}=
 (-a\,d''b\wedge d''c,\ b\,d'a\wedge d'c,
                                      c\,d'a\wedge d''b).
 \tag{NCP15}
\]
It suffices to verify this identity locally where all homogeneous
coordinates are nonzero, provided that the identity is established
as an actual finite native relation and is preserved under the
native unit substitutions used in that verification.  The previous
lemma then extends it across every proper-coordinate-zero point:
the formula (NCP14) was already actual there, and all of its scalar
arguments and curvatures coincide with the locally modified unit
presentation.  This is not an analytic-density argument.

For original data on $W$, the genuine original-$K$ map $[f]$ gives
$L|_W\simeq[f]^*\mathcal O(1)$.  Its normalized native data are
\[
 C_0=R-D\Phi_0,\qquad \theta_v=\Phi_v-\Phi_0,
                 \widehat\Psi_\nu=\Psi_\nu-\Phi_0.
 \tag{NCP16}
\]
Interpret substitution of (NCP16) into the supplied formula as an
explicit finite actual presentation, not as a claim about arbitrary
current pullback.  Specialize the extra unit to $\varpi\in K$.
The conclusion of the previous paragraph transfers (NCP15) to this
specified original actual formula.

To restore the metric and the originally chosen face offsets, use
the explicit GR metric comparison and (NCP5).  For
$\ell=r-3\ge1$ put
\[
 h_0=-\Phi_0\sum_{j=0}^{\ell-1}C_0^jR^{\ell-1-j},
 \qquad Dh_0=C_0^\ell-R^\ell.
 \tag{NCP17}
\]
Since $\partial_2E_{a,b,c}=0$, the full QHT identity gives
\[
 (C_0^\ell-R^\ell)E_{a,b,c}
             =\partial_3\mathcal K_{h_0}(E_{a,b,c}).
 \tag{NCP18}
\]
For $\ell=0$ this correction is zero.  If $T_{\rm met}$ denotes
the explicit GR quadruple with
$\mathcal Q_{\rm original}-\mathcal Q_{\rm normalized}
=\partial_3T_{\rm met}$, the resulting completely specified
conditional witness is
\[
 \mathcal Q_{\boldsymbol\varepsilon}
 =\partial_3\left(
   \mathscr V|_{T=\varpi}+T_{\rm met}
   +\sum_\nu\sigma_\nu\mathcal T^{\rm CT}_{\nu;\gamma,\varepsilon_\nu}
   +\mathcal K_{h_0}(E_{a,b,c})\right)
                         +R^{r-3}E_{a,b,c}.
 \tag{NCP19}
\]
The cap-comparison and metric-comparison order is fixed by using
the normalized common-$\gamma$ cube first, then the original metric
with common $\gamma$, then the original prescribed offsets.

The independently extending triple in (NCP19) is defined on $N$.
However, (NCP19) is conditional on the actual supplied quadruple
and full native identity (NCP15).  This module does not supply them,
does not claim that any distributional equality is automatically
an actual identity, and does not claim general finite-cover coherence.
An original-$U$ Green application still requires the literal local
path/overlap data and the exact higher compatibility conditions of
the earlier FG and RSP inputs.

\paragraph{Exact inputs and exclusions.}
CS: \path{d0c06d21c2289d292fb735d6966d456b9e05f1e4fab10f263a5007169d79cacf}.
CT: \path{0c10e28336f9c779cc3dc4a35393b1fc52d1f6d7d154a4ec466c13d7c6514e95}.
GR: \path{c47cbd11e549bf94430182888bccffce20cad14ddc70f2929cab7ba04f5d3997}.
QHT: \path{aa10132dc6e8217dfa4c896699c76eaed90ce4d632abc5e0147cb9e278c423e0}.
CU: \path{d38d0c77f879e3ca80962fa0fac2e56f6ea159fd59971a0bd810fbadc0a91471}.
All are in the tropical/actual-tests branch.  No external theorem,
extension of the base field, arbitrary-minimum closure, arbitrary
proper pushforward, or general actual kernel exactness is imported.
